% ─────────────────────────────────────────────────────────────────────────────
\subsection*{Overview}

The Load Balancer Model (LBM) from \textbf{CKW10~§7.2} is a coupled STDEVS
system that processes tasks generated by a stochastic source and dispatches
them to one of two parallel servers.  Each server has no queue: an arriving
task is dropped if the server is busy (\textit{loss system}).
The experiment validates that the Cadmium STDEVS engine reproduces the
analytical task-loss probabilities given by Erlang's loss formula for
$M/M/1/1$ queues.

% ─────────────────────────────────────────────────────────────────────────────
\subsection*{Model Topology}

\begin{center}
\begin{tikzpicture}[
  box/.style={draw, rectangle, minimum width=2.0cm, minimum height=0.9cm,
              align=center, font=\small},
  >=Stealth
]
  \node[box] (lg) at (0,0)     {LG};
  \node[box] (wb) at (3.5,0)   {WB};
  \node[box] (s1) at (7,  0.9) {S1};
  \node[box] (s2) at (7, -0.9) {S2};

  \draw[->] (lg) -- node[above,font=\scriptsize] {task ($\lambda$)} (wb);
  \draw[->] (wb) -- node[above,font=\scriptsize] {$\lambda_1$} (s1);
  \draw[->] (wb) -- node[below,font=\scriptsize] {$\lambda_2$} (s2);
  \draw[->] (s1) -- ++(1.5,0) node[right,font=\scriptsize] {$\lambda_1'$ (done)};
  \draw[->] (s2) -- ++(1.5,0) node[right,font=\scriptsize] {$\lambda_2'$ (done)};
\end{tikzpicture}
\end{center}

\paragraph{Internal couplings.}
$\mathrm{IC} = \{\,(LG.\texttt{out} \to WB.\texttt{inp}),\;
                    (WB.\texttt{out}_1 \to S1.\texttt{inp}),\;
                    (WB.\texttt{out}_2 \to S2.\texttt{inp})\}$

% ─────────────────────────────────────────────────────────────────────────────
\subsection*{Atomic Model Specifications}

All models are expressed in DEVS-RND form (Activity~(c), CKW10 Fig.~3)
and use $r \sim U(0,1)$ drawn from \texttt{std::mt19937}.

\paragraph{Load Generator (LG).}

\[
  M^{LG} = \langle X, Y, S, \delta_{\mathrm{int}}, \delta_{\mathrm{ext}},
              \lambda, ta \rangle
\]
\begin{align*}
  X &= \emptyset, \quad Y = \{(\mathit{task}_1,\;\texttt{out}_1)\} \\
  S &= \mathbb{R}^+_0 \quad (\sigma = \text{next inter-departure time}) \\
  \delta_{\mathrm{int}}(s,\,r) &= -\tfrac{1}{d_r}\ln r, \quad r \sim U(0,1) \\
  \delta_{\mathrm{ext}} &= \bot \\
  \lambda(s) &= (\mathit{task}_1,\;\texttt{out}_1) \\
  ta(s) &= s
\end{align*}

The departure process is a Poisson process with rate $\lambda = d_r$ (tasks
per time unit).  The initial $\sigma$ is set to the mean $1/d_r$.

\paragraph{Weighted Balancer (WB).}

\[
  M^{WB} = \langle X, Y, S, \delta_{\mathrm{int}}, \delta_{\mathrm{ext}},
               \lambda, ta \rangle
\]
\begin{align*}
  X &= T \times \{\texttt{inp}\}, \quad
  Y  = T \times \{\texttt{out}_1,\, \texttt{out}_2\} \\
  S &= T \times \{1,2\} \times \mathbb{R}^+_0 \quad (w,\, p,\, \sigma) \\
  \delta_{\mathrm{ext}}\!\bigl((w,p,\sigma),\,e,\,(x_v,x_p),\,r\bigr)
        &= \bigl(x_v,\;\tilde p,\;0\bigr), \quad
           \tilde p = \begin{cases}1 & r < b_f \\ 2 & \text{otherwise}\end{cases} \\
  \delta_{\mathrm{int}}\!\bigl((w,p,\sigma),\,r\bigr) &= (w,\,p,\,+\infty) \\
  \lambda(w,p,\sigma) &= (w,\,\texttt{out}_p) \\
  ta(w,p,\sigma) &= \sigma
\end{align*}

After receiving a task, $\sigma = 0$ so WB outputs and then returns to
passive ($\sigma = +\infty$) instantaneously.

\paragraph{Servers S1 and S2.}

Let $\mu_i = 1/s_{t_i}$ be the service rate of server $i$.

\[
  M^{S_n} = \langle X, Y, S, \delta_{\mathrm{int}}, \delta_{\mathrm{ext}},
                \lambda, ta \rangle
\]
\begin{align*}
  X &= T \times \{\texttt{inp}\}, \quad Y = T \times \{\texttt{out}\} \\
  S &= T \times \{false,\,true\} \times \mathbb{R}^+_0
       \quad (w,\,\mathit{busy},\,\sigma) \\
  \delta_{\mathrm{ext}}\!\bigl((w,\mathit{busy},\sigma),\,e,\,(x_v,x_p),\,r\bigr)
    &= \begin{cases}
         (x_v,\;true,\;-\tfrac{1}{\mu_n}\ln r) & \neg\mathit{busy} \\
         (w,\;true,\;\sigma - e)                & \mathit{busy}
       \end{cases} \\
  \delta_{\mathrm{int}}\!\bigl((w,\mathit{busy},\sigma),\,r\bigr)
    &= (w,\;false,\;+\infty) \\
  \lambda(w,\mathit{busy},\sigma) &= (w,\;\texttt{out}) \\
  ta(w,\mathit{busy},\sigma) &= \sigma
\end{align*}

If busy when a task arrives, the task is discarded (loss system).
Service times are i.i.d.\ $\mathrm{Exp}(\mu_n)$.

% ─────────────────────────────────────────────────────────────────────────────
\subsection*{Analytical Validation}

Each server is an $M/M/1/1$ loss system (\textit{Erlang B} with $m = 1$).
The traffic intensities and theoretical loss probabilities are:
\begin{align}
  \lambda_1 &= b_f \cdot d_r, & \lambda_2 &= (1 - b_f)\cdot d_r \notag \\
  \rho_i    &= \lambda_i / \mu_i \notag \\
  P^{\mathrm{loss}}_i &= \frac{\rho_i}{1 + \rho_i}
    \label{eq:erlang-b} \\
  P^{\mathrm{loss}} &= b_f\, P^{\mathrm{loss}}_1 + (1 - b_f)\,P^{\mathrm{loss}}_2
    \label{eq:total-loss} \\
  \lambda' &= \lambda\,(1 - P^{\mathrm{loss}})
    \label{eq:throughput}
\end{align}

% ─────────────────────────────────────────────────────────────────────────────
\subsection*{Test Scenario (TS1)}

\begin{center}
\begin{tabular}{lll}
  \hline
  Parameter & Symbol & Value \\
  \hline
  Mean departure rate & $d_r$ & $10$ tasks/s \\
  Mean service time S1 & $s_{t_1}$ & $0.2$~s ($\mu_1 = 5$) \\
  Mean service time S2 & $s_{t_2}$ & $0.2$~s ($\mu_2 = 5$) \\
  Balance factor range & $b_f$ & $[0,\,1]$ in steps of $0.1$ \\
  Simulation time & $T$ & $10\,000$~s \\
  Repetitions per point & $N$ & $15$ \\
  \hline
\end{tabular}
\end{center}

Simulation results (Section~\ref{sec:ckw10-lbm-results}) match the
theoretical curves from equations~\eqref{eq:erlang-b}--\eqref{eq:throughput}
across all $11$ operating points, for all $15$ repetitions, reproducing
Figure~5 of CKW10.
